\section{Conclusions and open problems}\label{sec:conclusion}

Small treewidth alone does not make sparse mixed-integer problems
tractable. Two quantitative properties of the objective make tree
decompositions effective for certified global optimization on a box: upper
bounds on the coordinate curvature, which turn finite dynamic programming into
a valid lower bound and its min-marginals into a safe filter, and quadratic
growth, which keeps the filtered grids at a size independent of the accuracy.
Under weighted growth, CT (Algorithm~\ref{alg:ct}) certifies accuracy
$2^{-q}$ in $f(p,\bar\kappa)(I+q+1)^5$ bit operations, and under point growth
exact output takes $f_1(p,\kappa)(I+1)^{C_1}$. The algorithms need no knowledge of
the growth constant, and their certificates are checked by recomputation in
exact arithmetic and are valid without any growth assumption. The lower bounds of
Section~\ref{sec:limits} show that, unless P${}={}$NP, neither parameter can
be dropped and the dependence on $\kappa$ cannot be polylogarithmic; under
ETH the dependence on $p$ is exponential even for $\kappa\le2$; and under
rETH the exponent of $\kappa$ cannot be $o(p)$, already for integer box
quadratics. Exact recourse reduces the curvature that has to be discretized
or removes a residual problem from discretization altogether.

The certificates may also be useful outside complete algorithms. A
corrected-grid bound on a subbox is a valid node bound for spatial
branch-and-bound, and the min-marginal filter is bound tightening whose
reductions are certified; whether either pays off inside a solver is
untested. A replayable certificate also allows a floating-point incumbent to
be certified: its exact value together with a grid certificate gives a
rigorous gap. This matters because floating-point solvers satisfy constraints
only within tolerances (Section~\ref{sec:computation}). The experiments show
accuracy-independent grid sizes on designed families and on unplanted random
instances, including those on which growth could not be certified. On two
binary QPLIB instances the method is plain dynamic programming: the one whose
decomposition has bag size 20 is solved exactly in one stage, and the other
has bag size 94, so the width of the available decomposition, not the
accuracy, is the binding limit.

Several questions remain open.
\begin{enumerate}[label=(\arabic*)]
\item \emph{Negative curvature.} Let $g$ be a point-growth constant,
$\nu=\max\{0,-\lambda_{\min}(\nabla^2F)\}$ and $\kappa_\nu=\max\{1,\nu/g\}$.
Can continuous box QP be solved to accuracy $2^{-q}$ in time
$f'(p,\kappa_\nu)\poly(I+q)$ for a computable function $f'$?
Theorem~\ref{thm:cv} and Corollary~\ref{cor:cv-nu} give such a bound, with
$p$ the bag size after elimination, when certified responses of private
convex blocks reduce every coordinate curvature to at most $C_0\nu$ for a
constant $C_0$.
Proposition~\ref{prop:cv-limit} shows that recourse alone cannot ensure this,
and Proposition~\ref{lim:prop:moments} shows that matching local moments of
any finite order is not enough. These obstructions suggest measuring
curvature only where near-optimal points can lie, or using inequalities that
couple both sides of a separator.
\item \emph{Constants in the exponent.} Under weighted growth, each bag
table of CT has at most
$(C\sqrt{\bar\kappa}\log(n+2))^p$ entries. The factor $p^p$ in
$f(p,\kappa)=c_0(c_1p\sqrt\kappa)^p\kappa(1+\log_2\kappa)^2$ comes only from
the bound $\lceil\log_2(n+2)\rceil^p\le2p^p(n+2)$
(inequality~\eqref{eq:logabsorb}), which moves the logarithmic factor into
the polynomial part; the value-oracle lower bound of
Proposition~\ref{lim:prop:oracle} has no such factor. Can the logarithmic
factor in the grid size, and with it the factor $p^p$, be avoided, and is the
exponent $p/2$ of $\kappa$ optimal for explicit inputs, as it is for value
oracles?
\item \emph{Coupling constraints.} Is there a feasible rounding scheme for
totally unimodular, or more general, sparse constraints that works with
graded grids? Proposition~\ref{prop:tu-misaligned} and
Example~\ref{ex:tu-sum} show that curvature-only corrections on misaligned
grids and on a common graded pattern give invalid bounds, so a different
correction is needed. A partial step would keep the aligned uniform mesh
only on coordinates that occur in a constraint and use graded grids on the
others, which should replace the factor $n_c^{p/2}$ of
Theorem~\ref{thm:tu-approx} by $n^{p_c/2}$, where $p_c$ is the largest number
of constrained continuous coordinates in a bag; we have not worked out this
combination.
\item \emph{Optimal sets.} Suppose that the optimal set is finite and that
$r$ bounds the number of optimal values per coordinate. Can exact output then
be obtained in $f'(p,\kappa_S,r)\poly(I)$ bit operations for a computable
function $f'$? For optimal sets with continuous components, the grids
produced by filtering can need a number of nodes that grows with the accuracy
(Proposition~\ref{lim:prop:setgrowth}); which compact descriptions beyond the
classes of Section~\ref{sec:optsets} can be computed with a decomposition?
\item \emph{Exact polynomial output.} For continuous polynomial objectives the
minimizer can be irrational. Certificates that keep correlations between
coordinates, rather than signs over coordinate rectangles, seem necessary
under weak complementarity (Proposition~\ref{prop:weakcompl}); no construction
of parameterized size is known.
\end{enumerate}
